\section{Task Modality and Benchmark (Carrier)}
\label{sec:benchmark}
% 【中文】\section{任务模态与基准（载体）}

\bench{} is a \emph{two-tier} benchmark of long-chain EDA design-flow tasks across \textbf{13 EDA domains}---from schematic/netlist to SPICE simulation, physical verification, and digital synthesis---each derived from \emph{real coding-agent traces} of IC design engineers in a production environment. We use \bench{} as a \emph{carrier} for falsifiable cognition claims, not as the contribution itself.
% 【中文】
% \bench{} 是一个\emph{两层}的长链 EDA 设计流任务基准，跨越 \textbf{13 个 EDA 领域}——从原理图/网表到
% SPICE 仿真、物理验证与数字综合——每个任务源自生产环境中 IC 设计工程师的\emph{真实编程智能体轨迹}。
% 我们将 \bench{} 用作可证伪认知主张的\emph{载体}，而非贡献本身。

\textbf{Two-tier structure.} The \emph{Full} set comprises \FullScale{} tasks (continuously growing as new traces are curated; the exact count is not reported in this paper) curated from real IC-design engineer traces. The \emph{Lite} set is a fixed \LiteSet{}-task subset filtered from Full by domain-coverage balance (preserving all 13 domains) and chain-length distribution (approximating Full's), used for general evaluation under bounded token/time budgets; all arm-level experiments in this paper run on Lite. A frozen \Neval{}-run snapshot sampled from the Full set provides the descriptive Full-set baseline (below).
% 【中文】
% \textbf{两层结构。} \emph{Full} 集含 \FullScale{} 个任务（随新 trace 整理而持续增长；本文不报告确切
% 计数），均整理自真实 IC 设计工程师轨迹。\emph{Lite} 集是从 Full 按领域覆盖平衡（保留全部 13 域）
% 与链长分布（近似 Full）筛选的固定 \LiteSet{} 任务子集，用于在有限 token/时间预算下的一般评测；
% 本文所有臂级实验在 Lite 上运行。一份从 Full 集采样的冻结 \Neval{}-run 快照提供描述性 Full-set
% 基线（见下）。

\textbf{Design-flow contract.} Each task requires the agent to produce \texttt{submit.zip} by orchestrating remote EDA tools through an MCP bridge, scored against $\geq 8$ executable assertions with a \texttt{PreToolUse} hook that prevents answer leakage (reference/ and check.py never enter the workdir). The contract deliberately does not specify language or host---the agent must discover the right tool chain. Detailed benchmark construction is in Appendix~\ref{sec:benchmark-details}.
% 【中文】
% \textbf{设计流契约。} 每个任务要求智能体通过 MCP 网桥编排远程 EDA 工具产出 \texttt{submit.zip}，
% 按 $\geq 8$ 条可执行断言评分，并用 \texttt{PreToolUse} 钩子防止答案泄漏（reference/ 与 check.py
% 永不进入工作目录）。该契约故意不规定语言或主机——智能体必须自行发现正确的工具链。详细基准构建见
% 附录~\ref{sec:benchmark-details}。

\textbf{Evaluation subset.} Arm-level experiments run on the Lite \LiteSet{}-task set; the \emph{minimal matrix} for arm comparison is the top-120 of Lite by chain length (power-constrained arm-comparison anchor), and the \emph{ideal matrix} is Lite $\times$ 3 repeats (\S\ref{sec:hypotheses}). A frozen 696-run Full-set snapshot (2026-11-09; DeepSeek-V4-Pro, Cline v3.0.51, \texttt{thinking=low}, 2400s, 8 workers) provides the descriptive Full-set baseline: \MzeroPass{} Pass@1 (62 PASS / 558 FAIL / 70 ERROR / 6 ABSENT). Because 696 runs are sampled from Full (not Lite) and the top-120 is harder than the Full set, the minimal-matrix comparison anchor for arms is the M0 Pass@1 on the \emph{same top-120-of-Lite subset} (\TBD, to be measured with the arms), expected to differ from 10.0\% (\S\ref{sec:hypotheses}).
% 【中文】
% \textbf{评测子集。} 臂级实验在 Lite \LiteSet{}-任务集上运行；臂对比的\emph{最小矩阵}为按链长取
% Lite 的 top-120（受功效约束的臂对比锚点），\emph{理想矩阵}为 Lite $\times$ 3 重复
%（\S\ref{sec:hypotheses}）。一份冻结的 696-run Full-set 快照（2026-11-09；DeepSeek-V4-Pro、
% Cline v3.0.51、\texttt{thinking=low}、2400s 超时、8 workers）提供描述性 Full-set 基线：
% \MzeroPass{} Pass@1（62 PASS / 558 FAIL / 70 ERROR / 6 ABSENT）。由于 696 run 采样自 Full
% （非 Lite）且 top-120 比全集更难，各臂最小矩阵的比较锚点为\emph{同 top-120-of-Lite 子集}上实测的
% M0 Pass@1（\TBD，与各臂一同测量），该值预期偏离 10.0\%（\S\ref{sec:hypotheses}）。

\textbf{Reproducibility boundary.} Due to the production-origin of the traces (real IC design engineer sessions on commercial EDA tools), no data, code, or executable assets are released; reproducibility is supported via the detailed protocol description in this paper and an \emph{NDA-gated review package} for qualified reviewers. The falsifiability of our claims does not depend on proprietary licenses: the protocol and criteria suffice for an independent party to re-derive an equivalent benchmark on open-source EDA tools (Yosys/OpenROAD/Magic) and test the same hypotheses (\S\ref{sec:reproducibility}).
% 【中文】
% \textbf{可复现性边界。} 由于轨迹源自生产环境（真实 IC 设计工程师在商用 EDA 工具上的会话），不发布
% 任何数据、代码或可执行资产；可复现性通过本论文的详细协议描述与面向合格评审者的\emph{NDA 受控评审包}
% 支持。我们主张的可证伪性不依赖于专有许可：协议与判据足以让独立方在开源 EDA 工具
%（Yosys/OpenROAD/Magic）上重新推导出等价基准并检验同样假设（\S\ref{sec:reproducibility}）。
